\subsection{Experienced pilots and beginners}\label{sec:levels}

Does the route explain the same share for both wing levels?

Not quite (Table~\ref{tab:isolated}, Figure~\ref{fig:isolated}). For the \IsoPiecesExp{} pieces of
experienced pilots $\Delta\Hloc(1000~\mathrm{s})$ falls from \IsoBaseDHlocExp{} to
\IsoDHlocExp{} \IsoDHlocExpBand, \IsoExplainedExp\% explained. For the \IsoPiecesBeg{} pieces of
beginners it falls from \IsoBaseDHlocBeg{} to \IsoDHlocBeg{} \IsoDHlocBegBand, \IsoExplainedBeg\%
explained. The beginners' baseline gap is larger, their share lower and their band wider, since they have
fewer pieces. Split by level, the gap that remains is not distinguished from zero.

\begin{figure}
\centering
\includegraphics[width=\textwidth]{../../out/d120_eh500_ez300_th45_g120_q0.1/open_declared/isolated.pdf}
\caption{Left: $\Hloc(\tau)$ of group time (orange) and of the same route flown alone (blue), by wing
level, with the number of pieces in the panel titles. Continuous: all the pieces, up to 1000~s; dashed:
only the pieces long enough for the longest lag, a subset of them, whose lags go further. Right: the difference group minus control of $\Hloc(1000~\mathrm{s})$ and of the
statistics of climbs and glides, with 5--95\% bands, by wing level.}
\label{fig:isolated}
\end{figure}

\paragraph{What we keep.} Both levels lose most of their gap on the same route, experienced pilots
more than beginners, and neither remaining gap is distinguished from zero on its own.
